%!TEX program = xelatex
% ============================================================
%  הרצאה 6 — עבודה ואנרגיה II
%  (Work & Energy II — Conservative Forces & Energy Conservation)
% ============================================================
\documentclass[12pt, a4paper]{article}

\usepackage{amsmath, amssymb, mathtools}
\usepackage{xcolor}
\usepackage{graphicx}

\usepackage{tcolorbox}
\tcbuselibrary{skins, breakable}
\tcbset{
  resultbox/.style={
    enhanced, colback=white, colframe=black,
    boxrule=0.8pt, arc=3pt,
    left=2pt, right=2pt, top=1.5pt, bottom=1.5pt,
  },
  notebox/.style={
    enhanced, breakable,
    colback=gray!8, colframe=gray!50,
    boxrule=0.8pt, arc=3pt,
    left=2pt, right=2pt, top=1.5pt, bottom=1.5pt,
    fonttitle=\bfseries,
  },
  warnbox/.style={
    enhanced, breakable,
    colback=red!5, colframe=red!60,
    boxrule=0.8pt, arc=3pt,
    left=3pt, right=3pt, top=2pt, bottom=2pt,
    fonttitle=\bfseries,
  },
  defbox/.style={
    enhanced, breakable,
    colback=blue!5, colframe=blue!50,
    boxrule=0.8pt, arc=3pt,
    left=3pt, right=3pt, top=2pt, bottom=2pt,
    fonttitle=\bfseries,
  },
}

\usepackage{tikz}
\usetikzlibrary{arrows.meta, calc, decorations.pathreplacing, patterns}
\usepackage{pgfplots}          % velocity/energy graphs (standard in TeXLive)
\pgfplotsset{compat=1.18}
\usepackage[a4paper, top=2cm, bottom=2cm,
            left=2cm, right=2cm]{geometry}
\usepackage{setspace}
\setstretch{1.3}

% ── listings must be loaded BEFORE polyglossia/bidi ──────────
\usepackage{listings}
\usepackage{titlesec}
\titleformat{\section}{\large\bfseries}{\thesection.}{0.5em}{}[\titlerule]
\titleformat{\subsection}{\bfseries}{\thesubsection}{0.5em}{}

\setlength{\headheight}{15pt}
\usepackage{fancyhdr}
\pagestyle{fancy}
\fancyhf{}
\rhead{\textenglish{\textbf{Work \& Energy II}}}
\lhead{\thepage}
\renewcommand{\headrulewidth}{0.3pt}

\usepackage{fontspec}
\usepackage{polyglossia}
\setmainlanguage{hebrew}
\setotherlanguage{english}
\setmainfont{Times New Roman}
\newfontfamily\hebrewfont{Times New Roman}
\newfontfamily\englishfont{Times New Roman}
\newfontfamily\hebrewfonttt{Courier New}
\setmonofont{Courier New}      % monospace for Python listings

% ── Mandatory bidi + TikZ fix (after polyglossia) ───────────
% Forces every tikzpicture into LTR so Hebrew labels do not throw
% pgfkeys '/tikz/{right}' errors or render mirrored.
\usepackage{etoolbox}
\BeforeBeginEnvironment{tikzpicture}{\begin{LTR}}
\AfterEndEnvironment{tikzpicture}{\end{LTR}}

% ── Python code listings (package loaded earlier, before bidi) ─
\definecolor{codekw}{RGB}{0,0,180}
\definecolor{codecmt}{RGB}{34,120,34}
\definecolor{codestr}{RGB}{170,80,0}
\lstdefinestyle{pycode}{
  language=Python,
  basicstyle=\ttfamily\footnotesize,
  keywordstyle=\color{codekw}\bfseries,
  commentstyle=\color{codecmt}\itshape,
  stringstyle=\color{codestr},
  numbers=left,
  numberstyle=\tiny\color{gray},
  numbersep=8pt,
  showstringspaces=false,
  breaklines=true,
  frame=single,
  rulecolor=\color{gray!50},
  framerule=0.6pt,
  framesep=5pt,
  columns=fullflexible,
  keepspaces=true,
  aboveskip=6pt, belowskip=4pt,
}

\newcommand{\hvec}[1]{\overrightarrow{#1}}

% ============================================================
\begin{document}
% ============================================================

% ── Title ────────────────────────────────────────────────────
\begin{center}
  {\LARGE\bfseries עבודה ואנרגיה \textenglish{(Work \& Energy II)}}\\[6pt]
  \rule{\linewidth}{0.8pt}
\end{center}
\vspace{0.5cm}

בהרצאה הקודמת הגדרנו עבודה כאינטגרל מסלולי של הכוח, וראינו את משפט עבודה--אנרגיה. בהרצאה זו נגלה ש\textbf{העבודה תלויה במסלול} עבור כוחות מסוימים אך \textbf{אינה תלויה במסלול} עבור אחרים. הבחנה זו תוביל אותנו אל מושג \textbf{הכוח המשמר}, אל \textbf{האנרגיה הפוטנציאלית}, ואל אחד העקרונות החשובים בפיזיקה --- \textbf{שימור האנרגיה המכנית}.

% ============================================================
\section*{0. חזרה --- עבודה, אנרגיה והספק}
% ============================================================

לפני שנמשיך, נסכם את הכלים שהוגדרו בהרצאה הקודמת.

\begin{tcolorbox}[notebox, title={תקציר ההרצאה הקודמת}]
\begin{itemize}
  \item \textbf{אנרגיה קינטית:} \quad \textenglish{$E_k = \tfrac{1}{2} m v^2$}.
  \item \textbf{עבודה (הגדרה):} אינטגרל מסלולי של הכוח,
        \[
          W = \int_{\hvec{r}_1}^{\hvec{r}_2} \hvec{F}\cdot d\hvec{r}
            = \int \bigl(F_x\,dx + F_y\,dy + F_z\,dz\bigr).
        \]
  \item \textbf{הספק:} \quad
        \textenglish{$W = \displaystyle\int_{t_1}^{t_2} P\,dt
                          = \int \hvec{F}\cdot\hvec{v}\,dt
                          = \int \hvec{F}\cdot d\hvec{r}$}.
  \item \textbf{משפט עבודה--אנרגיה:} עבודת הכוח השקול שווה לשינוי האנרגיה הקינטית,
        \[
          W_{\text{שקול}} = \Delta E_k = E_{k,2} - E_{k,1}.
        \]
\end{itemize}
\end{tcolorbox}

% ============================================================
\section*{1. תלות העבודה במסלול --- דוגמה מרכזית}
% ============================================================

\subsection*{1.1 הצגת הבעיה}

נתון כוח התלוי במיקום:
\begin{tcolorbox}[resultbox]
\[
  \hvec{F}(x,y,z) = \alpha\,x^2 y\,\hat{x} + \beta\,x\,\hat{y},
  \qquad
  \alpha = 1\;\tfrac{\text{N}}{\text{m}^3},
  \quad
  \beta  = 1\;\tfrac{\text{N}}{\text{m}}.
\]
\end{tcolorbox}

מבקשים לחשב את עבודת הכוח בין הנקודות
\[
  \hvec{r}_1 = \hvec{0},
  \qquad
  \hvec{r}_2 = 1\cdot\hat{x} + 1\cdot\hat{y}\;\;[\text{m}].
\]
על פי ההגדרה,
\[
  W = \int_{\hvec{r}_1}^{\hvec{r}_2} \hvec{F}\cdot d\hvec{r}
    = \int_0^1 \alpha\,x^2 y\,dx + \int_0^1 \beta\,x\,dy .
\]

\begin{tcolorbox}[warnbox, title={מה עושים עם זה?}]
\textbf{ללא ידיעת המסלול לא ניתן להמשיך בחישוב.} באינטגרל הראשון (לפי \textenglish{$x$}) מופיעה הקואורדינטה \textenglish{$y$}, ובשני (לפי \textenglish{$y$}) מופיעה \textenglish{$x$}. כדי לבצע את האינטגרציה חייבים \textbf{משוואת מסלול} --- ביטוי של \textenglish{$y$} כתלות ב-\textenglish{$x$} (ולהפך) --- ולהציבה.
\end{tcolorbox}

נדגים שני מסלולים שונים המחברים בין אותן שתי נקודות:
\begin{enumerate}
  \item הקו הישר \textenglish{$y = x$}.
  \item הפרבולה \textenglish{$y = x^2$}.
\end{enumerate}

\begin{center}
\begin{tikzpicture}[scale=3.0, >={Stealth[length=2.5mm]}]
  % axes
  \draw[->, thick] (-0.08,0) -- (1.22,0) node[right] {\textenglish{$x$}};
  \draw[->, thick] (0,-0.08) -- (0,1.22) node[above] {\textenglish{$y$}};
  % endpoints
  \fill (0,0)  circle (0.6pt) node[below left]  {\textenglish{$\hvec{r}_1$}};
  \fill (1,1)  circle (0.6pt) node[above right] {\textenglish{$\hvec{r}_2$}};
  % straight line y = x
  \draw[blue, very thick] (0,0) -- (1,1);
  \node[blue] at (0.40,0.56) {\textenglish{$y=x$}};
  % parabola y = x^2
  \draw[red, very thick] plot[domain=0:1, samples=80] (\x,{\x*\x});
  \node[red] at (0.86,0.55) {\textenglish{$y=x^2$}};
\end{tikzpicture}
\end{center}

\subsection*{1.2 מסלול 1 --- הקו הישר \textenglish{$y=x$}}

בקו הישר \textenglish{$y=x$}: באינטגרל לפי \textenglish{$x$} נציב \textenglish{$y=x$}, ובאינטגרל לפי \textenglish{$y$} נציב \textenglish{$x=y$}:
\[
  W_1 = \int_0^1 x^2\,(x)\,dx + \int_0^1 (y)\,dy
      = \int_0^1 x^3\,dx + \int_0^1 y\,dy
      = \frac{1}{4} + \frac{1}{2}.
\]
\begin{tcolorbox}[resultbox]
\[
  W_1 = \frac{3}{4}\;\text{J}
\]
\end{tcolorbox}

\subsection*{1.3 מסלול 2 --- הפרבולה \textenglish{$y=x^2$}}

בפרבולה \textenglish{$y=x^2$}: באינטגרל לפי \textenglish{$x$} נציב \textenglish{$y=x^2$}, ובאינטגרל לפי \textenglish{$y$} נציב \textenglish{$x=\sqrt{y}=y^{1/2}$}:
\[
  W_2 = \int_0^1 x^2\,(x^2)\,dx + \int_0^1 (y^{1/2})\,dy
      = \int_0^1 x^4\,dx + \int_0^1 y^{1/2}\,dy
      = \frac{1}{5} + \frac{2}{3}.
\]
\begin{tcolorbox}[resultbox]
\[
  W_2 = \frac{1}{5} + \frac{2}{3} = \frac{13}{15}\;\text{J}
\]
\end{tcolorbox}

\subsection*{1.4 מסקנה}

\[
  W_1 = \frac{3}{4} = \frac{45}{60},
  \qquad
  W_2 = \frac{13}{15} = \frac{52}{60}
  \qquad\Longrightarrow\qquad
  W_1 \ne W_2 .
\]

\begin{tcolorbox}[defbox, title={תלות במסלול}]
עבור כוח זה \textbf{העבודה תלויה במסלול}: אותן נקודות קצה נתנו עבודות שונות. זוהי תכונה של חלק מהכוחות --- ומכאן ההבחנה היסודית:
\begin{itemize}
  \item יש כוחות שעבורם העבודה \textbf{אינה תלויה במסלול} --- \textbf{כוחות משמרים} (\textenglish{conservative}).
  \item יש כוחות שעבורם העבודה \textbf{תלויה במסלול} --- \textbf{כוחות אי-משמרים} (\textenglish{non-conservative}).
\end{itemize}
\end{tcolorbox}

% ============================================================
\section*{2. כוחות משמרים --- הגדרה ותכונה}
% ============================================================

\subsection*{2.1 הגדרה}

\begin{tcolorbox}[defbox, title={כוח משמר (\textenglish{Conservative Force})}]
\textbf{כוח משמר} הוא כוח שעבודתו \textbf{אינה תלויה במסלול} שבין שתי נקודות, אלא רק בנקודות הקצה.
\end{tcolorbox}

\subsection*{2.2 תכונה: עבודה על מסלול סגור}

\begin{tcolorbox}[resultbox, title={תכונת הכוח המשמר}]
\[
  \oint \hvec{F}\cdot d\hvec{r} = 0
  \qquad\text{(כוח משמר)}
\]
העבודה של כוח משמר על \textbf{מסלול סגור} שווה לאפס.
\end{tcolorbox}

\textbf{הוכחה.} נתבונן בשתי נקודות \textenglish{A} ו-\textenglish{B} ובשני מסלולים המחברים ביניהן: מסלול 1 (עבודה \textenglish{$W_1$}) ומסלול 2 (עבודה \textenglish{$W_2$}).

\begin{center}
\begin{tikzpicture}[scale=1.3, >={Stealth[length=3mm]}]
  % endpoints
  \fill (0,0)   circle (2pt) node[below left]  {\textenglish{A}};
  \fill (3,1.6) circle (2pt) node[above right] {\textenglish{B}};
  % path 1 (upper arc) A -> B
  \draw[blue, very thick, ->] (0,0) .. controls (0.5,1.7) and (1.7,2.05) .. (3,1.6);
  \node[blue, above] at (1.35,1.95) {\textenglish{$W_1$}};
  % path 2 (lower arc) A -> B
  \draw[red, very thick, ->] (0,0) .. controls (1.3,-0.45) and (2.5,0.1) .. (3,1.6);
  \node[red, below] at (1.6,-0.32) {\textenglish{$W_2$}};
\end{tikzpicture}
\end{center}

מאחר שהכוח משמר, העבודה אינה תלויה במסלול, ולכן
\[
  W_1 = W_2 .
\]
כעת נבנה \textbf{מסלול סגור}: נלך מ-\textenglish{A} ל-\textenglish{B} דרך מסלול 1, ונחזור מ-\textenglish{B} ל-\textenglish{A} דרך מסלול 2 \textbf{בכיוון ההפוך}. עבודת החזרה היא \textenglish{$W_3 = -W_2$} (היפוך כיוון התנועה הופך את סימן העבודה). לכן העבודה על המסלול הסגור:
\[
  W_{\text{סגור}} = W_1 + W_3 = W_1 - W_2 = 0 .
\]
\hfill$\blacksquare$

% ============================================================
\section*{3. אנרגיה פוטנציאלית ומשפט עבודה--אנרגיה II}
% ============================================================

\subsection*{3.1 פירוק הכוח השקול}

נמיין את הכוח השקול לרכיב משמר ולרכיב אי-משמר:
\begin{tcolorbox}[resultbox]
\[
  \hvec{F} = \hvec{F}_C + \hvec{F}_{NC}
\]
\end{tcolorbox}
כאשר \textenglish{$\hvec{F}_C$} הוא סך הכוחות המשמרים ו-\textenglish{$\hvec{F}_{NC}$} סך הכוחות האי-משמרים.

נציב במשפט עבודה--אנרגיה (כרגיל):
\[
  W = \int_{\hvec{r}_1}^{\hvec{r}_2}\!\bigl(\hvec{F}_C + \hvec{F}_{NC}\bigr)\cdot d\hvec{r}
    = K_2 - K_1 ,
\]
ולכן
\[
  \int_{\hvec{r}_1}^{\hvec{r}_2}\!\hvec{F}_{NC}\cdot d\hvec{r}
   = (K_2 - K_1) - \int_{\hvec{r}_1}^{\hvec{r}_2}\!\hvec{F}_C\cdot d\hvec{r} .
\]

\subsection*{3.2 הגדרת האנרגיה הפוטנציאלית}

נגדיר את האנרגיה הפוטנציאלית של הרכיב המשמר באמצעות העבודה שלו:

\begin{tcolorbox}[defbox, title={אנרגיה פוטנציאלית (\textenglish{Potential Energy})}]
\[
  U(\hvec{r}_2) - U(\hvec{r}_1)
   = -\int_{\hvec{r}_1}^{\hvec{r}_2}\!\hvec{F}_C\cdot d\hvec{r}
\]
\end{tcolorbox}

\begin{tcolorbox}[notebox, title={הערה}]
ההגדרה תקפה \textbf{רק עבור כוח משמר}: היות שהעבודה אינה תלויה במסלול, הביטוי \textenglish{$U(\hvec{r})$} מוגדר היטב כפונקציה של המיקום בלבד. עבור כוח אי-משמר אין אנרגיה פוטנציאלית. כמו כן \textenglish{$U$} מוגדרת עד כדי \textbf{קבוע} (בחירת נקודת הייחוס), שכן רק \textbf{הפרשי} פוטנציאל הם בעלי משמעות פיזיקלית.
\end{tcolorbox}

\subsection*{3.3 משפט עבודה--אנרגיה II ושימור האנרגיה}

נציב את הגדרת \textenglish{$U$} (כך ש-\textenglish{$-\int \hvec{F}_C\cdot d\hvec{r} = U_2 - U_1$}):
\[
  \int_{\hvec{r}_1}^{\hvec{r}_2}\!\hvec{F}_{NC}\cdot d\hvec{r}
   = (K_2 - K_1) + (U_2 - U_1)
   = (K_2 + U_2) - (K_1 + U_1).
\]
נגדיר את \textbf{האנרגיה המכנית}:
\begin{tcolorbox}[resultbox, title={אנרגיה מכנית}]
\[
  E = K + U
\]
\end{tcolorbox}

ומתקבל משפט עבודה--אנרגיה בנוסחה השנייה:

\begin{tcolorbox}[resultbox, title={משפט עבודה--אנרגיה II}]
\[
  W_{NC} = \int_{\hvec{r}_1}^{\hvec{r}_2}\!\hvec{F}_{NC}\cdot d\hvec{r}
         = E_2 - E_1 = \Delta E
\]
עבודת הכוחות \textbf{האי-משמרים} שווה לשינוי האנרגיה המכנית.
\end{tcolorbox}

\begin{tcolorbox}[warnbox, title={מקרה פרטי חשוב ביותר --- שימור האנרגיה}]
אם \textbf{אין כוחות אי-משמרים} (או שאינם עושים עבודה), אזי \textenglish{$W_{NC} = 0$}, ולכן
\[
  \Delta E = 0
  \qquad\Longrightarrow\qquad
  E = K + U = \text{const}.
\]
\textbf{האנרגיה המכנית נשמרת.} זו בדיוק הסיבה לשם ``כוח משמר'': כוח משמר \emph{שומר} על האנרגיה המכנית --- כאשר האנרגיה הקינטית יורדת, הפוטנציאלית עולה באותו השיעור (והפוך).
\end{tcolorbox}

% ============================================================
\section*{4. סיווג הכוחות המוכרים}
% ============================================================

נסווג את הכוחות המוכרים לנו לשלוש קבוצות: \textbf{משמרים}, \textbf{אי-משמרים}, ו\textbf{כאלה שאינם עושים עבודה}. עבור כל כוח משמר נחשב את האנרגיה הפוטנציאלית שלו, לפי
\[
  E_p^{(F)}(\hvec{r}_2) - E_p^{(F)}(\hvec{r}_1)
   = -\int_{\hvec{r}_1}^{\hvec{r}_2}\!\hvec{F}\cdot d\hvec{r}
   = -W_{\hvec{r}_1\to\hvec{r}_2}
   =  W_{\hvec{r}_2\to\hvec{r}_1}.
\]

\begin{center}
\begin{tabular}{|c|c|c|}
\hline
\textbf{כוח} & \textbf{סיווג} & \textbf{הערה} \\
\hline
כוח הכובד          & משמר      & \textenglish{$E_p = mgh$} \\
\hline
מתיחות (חוט)        & אינו עושה עבודה & על המערכת / כשגוף נייח \\
\hline
נורמל              & אינו עושה עבודה & ניצב למהירות \\
\hline
חיכוך סטטי          & אינו עושה עבודה & נקודת המגע במנוחה \\
\hline
חיכוך קינטי         & אי-משמר   & תלוי במסלול \\
\hline
התנגדות אוויר (גרר) & אי-משמר   & תלוי במסלול \\
\hline
\end{tabular}
\end{center}

\subsection*{4.1 כוח הכובד --- כוח משמר}

\begin{tcolorbox}[notebox, title={למה כל כוח קבוע הוא משמר?}]
כוח הכובד \textenglish{$\hvec{F} = -mg\,\hat{y}$} הוא \textbf{כוח קבוע}, וכל כוח קבוע הוא משמר: רכיב \textenglish{$x$} של הכוח אינו תלוי ב-\textenglish{$y$} או ב-\textenglish{$z$}, וכך גם יתר הרכיבים --- ולכן העבודה אינה תלויה במסלול.
\end{tcolorbox}

נחשב את האנרגיה הפוטנציאלית הכבידתית. נבחר ציר \textenglish{$\hat{y}$} כלפי מעלה ונקודת ייחוס \textenglish{$E_p(\hvec{r}_0)=0$}:
\[
  E_p(\hvec{r}) = E_p(\hvec{r}_0) - \int_{\hvec{r}_0}^{\hvec{r}}\!\hvec{F}\cdot d\hvec{r}
   = 0 - \int_{y_0}^{y} (-mg)\,dy
   = mg\,(y - y_0).
\]
\begin{tcolorbox}[resultbox, title={אנרגיה פוטנציאלית כבידתית}]
\[
  E_p = mg\,h
  \qquad (h = y - y_0)
\]
\end{tcolorbox}

\subsection*{4.2 כוחות שאינם עושים עבודה}

\begin{tcolorbox}[notebox, title={מתיחות, נורמל וחיכוך סטטי}]
\begin{itemize}
  \item \textbf{מתיחות החוט} --- בחוט אידאלי הקושר שני גופים, המתיחות פועלת בכיוונים מנוגדים על שני קצותיו ובאותה היסט:
        \[
          W_T = T\,\Delta x + (-T)\,\Delta x = 0 .
        \]
  \item \textbf{נורמל} --- ניצב לכיוון התנועה, ולכן \textenglish{$\hvec{N}\cdot d\hvec{r}=0$}:
        \[
          W_N = N\,\Delta x + (-N)\,\Delta x = 0 .
        \]
  \item \textbf{חיכוך סטטי} --- בגלגול ללא החלקה נקודת המגע נמצאת \textbf{במנוחה רגעית} (\textenglish{$v=0$}), ולכן ההספק שלו מתאפס:
        \[
          P = f_s \cdot v = 0
          \quad\Longrightarrow\quad
          W_{f_s} = 0 .
        \]
\end{itemize}
\end{tcolorbox}

\textbf{חיכוך קינטי והתנגדות אוויר} הם כוחות \textbf{אי-משמרים}: כיוונם תמיד מנוגד למהירות, ולכן על מסלול סגור העבודה אינה מתאפסת (היא תמיד שלילית) --- ``בורחת'' מהמערכת כחום. עבורם \textbf{אין} אנרגיה פוטנציאלית.

% ============================================================
\section*{5. דוגמה --- מהירות מקסימלית של מטוטלת}
% ============================================================

\subsection*{5.1 הצגת הבעיה}

משחררים מטוטלת באורך \textenglish{$L$} \textbf{ממנוחה} מזווית \textenglish{$\theta_0$} מהאנך. בהזנחת התנגדות האוויר --- מהי המהירות \textbf{המקסימלית} שתשיג המטוטלת במהלך תנועתה?

\begin{center}
\begin{tikzpicture}[scale=1.35, >={Stealth[length=2.5mm]}]
  % ceiling
  \fill[pattern=north east lines, pattern color=gray!70] (-0.7,2.4) rectangle (0.7,2.62);
  \draw[thick] (-0.7,2.4) -- (0.7,2.4);
  \fill (0,2.4) circle (1.2pt);
  % vertical dashed
  \draw[gray!70, dashed] (0,2.4) -- (0,0.30);
  % string at theta0 = 35 deg, L = 2.0  -> bob (1.147,0.762)
  \draw[thick] (0,2.4) -- (1.147,0.762);
  \node at (0.46,1.66) {\textenglish{$L$}};
  % bob (released position)
  \shade[ball color=black] (1.147,0.762) circle (5pt);
  % bob at lowest point (faint)
  \draw[fill=gray!30] (0,0.40) circle (5pt);
  % angle arc (downward vertical -> string)
  \draw[->] (0,1.55) arc (270:305:0.85);
  \node at (0.33,1.50) {\textenglish{$\theta_0$}};
  % height h
  \draw[gray, dotted] (1.147,0.762) -- (1.80,0.762);
  \draw[gray, dotted] (0,0.40) -- (1.80,0.40);
  \draw[<->, gray] (1.62,0.762) -- (1.62,0.40);
  \node[gray, right] at (1.66,0.58) {\textenglish{$h = L(1-\cos\theta_0)$}};
\end{tikzpicture}
\end{center}

\subsection*{5.2 הדרך הקשה --- משוואת התנועה}

\begin{tcolorbox}[warnbox, title={למה לא לפתור ישירות?}]
שימוש בחוק השני של ניוטון בכיוון המשיק נותן את משוואת התנועה
\[
  -mg\sin\theta = mL\,\ddot{\theta}.
\]
זוהי משוואה דיפרנציאלית \textbf{לא ליניארית מסדר שני} (בגלל \textenglish{$\sin\theta$}). אין לנו פתרון אנליטי נוח, ולכן \textbf{לא נוכל} להמשיך בדרך זו בקלות.
\end{tcolorbox}

\subsection*{5.3 הדרך הקלה --- שימור אנרגיה}

הכוחות הפועלים על המטוטלת הם:
\begin{itemize}
  \item \textbf{כוח הכובד} --- כוח משמר.
  \item \textbf{מתיחות החוט} --- ניצבת תמיד למהירות, \textbf{אינה עושה עבודה}.
\end{itemize}
מאחר שאין כוח אי-משמר העושה עבודה --- \textbf{האנרגיה המכנית נשמרת}. האנרגיה הפוטנציאלית מינימלית בנקודה הנמוכה ביותר, ולכן שם תתקבל המהירות המקסימלית. נבחר \textenglish{$E_p = 0$} בנקודה הנמוכה ביותר במסלול.

\textbf{אנרגיה התחלתית} (במנוחה בזווית \textenglish{$\theta_0$}, בגובה \textenglish{$h = L(1-\cos\theta_0)$}):
\[
  E_i = E_k + E_p = 0 + mg\,L\,(1 - \cos\theta_0).
\]
\textbf{אנרגיה בנקודה הנמוכה ביותר} (\textenglish{$E_p = 0$}, מהירות מקסימלית):
\[
  E_f = E_k + E_p = \tfrac{1}{2} m v^2 + 0.
\]
משימור האנרגיה \textenglish{$E_i = E_f$}:
\[
  \tfrac{1}{2} m v^2 = mg\,L\,(1 - \cos\theta_0).
\]
\begin{tcolorbox}[resultbox, title={מהירות מקסימלית של מטוטלת}]
\[
  v_{\max} = \sqrt{\,2 g L\,(1 - \cos\theta_0)\,}
\]
\end{tcolorbox}

\begin{tcolorbox}[notebox, title={המסר}]
מה ש``נתקע'' בפתרון ישיר של משוואת התנועה נפתר \textbf{בשורה אחת} בעזרת שימור אנרגיה. זהו כוחה של גישת האנרגיה: כאשר רוצים גודל \emph{סקלרי} (מהירות בנקודה) ואין צורך בתלות המלאה בזמן --- שימור האנרגיה לרוב הדרך הקצרה.
\end{tcolorbox}

\subsection*{5.4 החלפת אנרגיה לאורך התנועה}

הגרף הבא מראה כיצד האנרגיה הקינטית והפוטנציאלית מתחלפות זו בזו לאורך הנדנוד, בעוד \textbf{הסכום} (האנרגיה המכנית) נשאר קבוע. הציר האנכי במונחי \textenglish{$mgL$}, והדוגמה עבור \textenglish{$\theta_0 = 1\,\text{rad}$}:

\begin{center}
\begin{tikzpicture}
\begin{axis}[
  width=11cm, height=6cm,
  xlabel={\textenglish{$\theta$ [rad]}},
  ylabel={\textenglish{Energy $[\,mgL\,]$}},
  domain=-1:1, samples=120,
  axis lines=left,
  xmin=-1.05, xmax=1.05, ymin=0, ymax=0.55,
  grid=both, grid style={gray!20},
  legend pos=south east,
]
\addplot[blue, very thick]        {1 - cos(deg(x))};
\addlegendentry{\textenglish{$E_p = mgL(1-\cos\theta)$}}
\addplot[red, very thick]         {cos(deg(x)) - 0.5403};
\addlegendentry{\textenglish{$E_k$}}
\addplot[black, dashed, thick]    {0.4597};
\addlegendentry{\textenglish{$E = E_k + E_p$ (const)}}
\end{axis}
\end{tikzpicture}
\end{center}

% ============================================================
\section*{6. דוגמה --- מהירות מקסימלית של מכונית}
% ============================================================

\subsection*{6.1 הצגת הבעיה}

מכונית שההספק המקסימלי של מנועה הוא \textenglish{$P$}, נתונה להשפעת התנגדות אוויר מן הצורה \textenglish{$\hvec{f} = -c v^2\,\hat{v}$}. מהי מהירותה המקסימלית, ומהי צריכת הדלק לק``מ במהירות זו?

\begin{tcolorbox}[notebox, title={נתונים}]
\begin{itemize}
  \item הספק מקסימלי: \textenglish{$P = 100\;\text{hp} \approx 74.6\;\text{kW}$}.
  \item מקדם גרר: \textenglish{$c = 0.5\;\tfrac{\text{kg}}{\text{m}}$}.
  \item תכולת אנרגיה של דלק: \textenglish{$40\;\text{MJ}/\ell$} (מגה-ג'אול לליטר).
  \item ניצולת המנוע: \textbf{20\%} (יתר האנרגיה הופכת לחום).
\end{itemize}
\end{tcolorbox}

\subsection*{6.2 דיאגרמת כוחות}

\begin{center}
\begin{tikzpicture}[scale=1.0, >={Stealth[length=3mm]}]
  % ground
  \fill[pattern=north east lines, pattern color=gray!60] (-0.6,-0.2) rectangle (5.4,0);
  \draw[thick] (-0.6,0) -- (5.4,0);
  % car body
  \draw[thick, fill=red!25, rounded corners=3pt] (1.4,0.25) rectangle (3.4,0.9);
  \draw[thick, fill=red!25, rounded corners=6pt] (1.9,0.9) rectangle (2.9,1.3);
  \fill[black] (1.9,0.25) circle (4.5pt);
  \fill[black] (2.9,0.25) circle (4.5pt);
  % forces
  \draw[->, very thick, blue]        (2.4,1.3)  -- (2.4,2.35) node[above] {\textenglish{$N$}};
  \draw[->, very thick, blue]        (2.4,0.25) -- (2.4,-1.0) node[below] {\textenglish{$mg$}};
  \draw[->, very thick, green!55!black] (3.4,0.6) -- (4.8,0.6) node[right] {\textenglish{$f_s$ (drive)}};
  \draw[->, very thick, red]         (1.4,0.6)  -- (0.0,0.6) node[left]  {\textenglish{$c v^2$ (drag)}};
\end{tikzpicture}
\end{center}

הכוחות: כוח הכובד, נורמל וחיכוך סטטי \textbf{אינם עושים עבודה}. כוח התנגדות האוויר \textbf{עושה עבודה}, וההספק שלו \textenglish{$P_{\text{drag}} = \hvec{f}\cdot\hvec{v} = -c v^2 \cdot v = -c v^3$}. בנוסף, הכוח הפנימי של המנוע עושה עבודה בהספק מקסימלי נתון \textenglish{$P$}.

\subsection*{6.3 פתרון --- מאזן הספקים}

במהירות המקסימלית התאוצה אפס, ולכן אין שינוי באנרגיה הקינטית --- עבודת הכוח השקול אפס, וכך גם סך ההספקים:
\[
  P_{\text{net}} = P - c v^3 = 0 .
\]
\begin{tcolorbox}[resultbox, title={מהירות מקסימלית}]
\[
  v_{\max} = \sqrt[3]{\dfrac{P}{c}}
\]
\end{tcolorbox}

הצבת הנתונים:
\[
  v_{\max} = \sqrt[3]{\frac{74600}{0.5}}
           = \sqrt[3]{1.49\times10^{5}}
           \approx 53\;\tfrac{\text{m}}{\text{s}}
           \;\approx\; 50\;\tfrac{\text{m}}{\text{s}}\ \text{(מעוגל)}.
\]

\subsection*{6.4 צריכת דלק לק``מ}

לנסיעה של \textenglish{$1000\;\text{m}$} במהירות \textenglish{$50\;\tfrac{\text{m}}{\text{s}}$} דרוש זמן
\[
  t = \frac{1000}{50} = 20\;\text{s}.
\]
העבודה שעושה המנוע בזמן זה:
\[
  W = P\cdot t = 74600 \times 20 \approx 1.5\;\text{MJ}.
\]
האנרגיה הזמינה בליטר דלק (לאחר ניצולת \textenglish{$20\%$}):
\[
  40\;\text{MJ}/\ell \times 0.20 = 8\;\text{MJ}/\ell .
\]
לכן צריכת הדלק לק``מ:
\[
  \frac{1.5\;\text{MJ}}{8\;\text{MJ}/\ell} \approx 0.2\;\ell
  \qquad\Longrightarrow\qquad
  \approx 5\;\tfrac{\text{km}}{\ell}.
\]

\begin{tcolorbox}[warnbox, title={המסקנה}]
\textenglish{$\approx 5\;\text{km}/\ell$} --- \textbf{יקר!} הצריכה גדלה כ-\textenglish{$v^3$} (ההספק נגד הגרר), ולכן \textbf{עדיף לנסוע הרבה יותר לאט}: הורדת המהירות חוסכת דלק בקצב חד.
\end{tcolorbox}

\subsection*{6.5 אימות נומרי בפייתון}

\noindent\textbf{חישוב המהירות המקסימלית וצריכת הדלק:}

\begin{LTR}
\begin{lstlisting}[style=pycode]
# Maximum speed and fuel economy from a power balance: P = c v^3
P  = 100 * 745.7    # engine power [W]  (100 hp)
c  = 0.5            # drag coefficient [kg/m]
E_fuel = 40e6       # energy content of fuel [J/liter]
eff = 0.20          # engine efficiency

v_max = (P / c) ** (1/3)            # P - c v^3 = 0  ->  v = (P/c)^(1/3)
print(f"v_max = {v_max:.1f} m/s")   # ~ 53 m/s

# Fuel for 1 km at v_max
d = 1000.0                          # distance [m]
t = d / v_max                       # travel time [s]
W = P * t                           # engine work over the km [J]
E_per_liter = E_fuel * eff          # usable energy per liter [J]
liters = W / E_per_liter            # liters for 1 km
print(f"t = {t:.1f} s,  W = {W/1e6:.2f} MJ")
print(f"fuel = {liters:.2f} L/km  ->  {1/liters:.1f} km/L")
\end{lstlisting}
\end{LTR}

% ============================================================
\section*{סיכום: מבט-על}
% ============================================================

\begin{tcolorbox}[warnbox, title={נקודות מפתח לזכור}]
\begin{enumerate}
  \item \textbf{תלות במסלול:} העבודה של חלק מהכוחות תלויה במסלול. דוגמה: \textenglish{$\hvec{F}=\alpha x^2 y\hat{x}+\beta x\hat{y}$} נתן \textenglish{$W_1=\tfrac{3}{4}$} לעומת \textenglish{$W_2=\tfrac{13}{15}$}.
  \item \textbf{כוח משמר:} העבודה אינה תלויה במסלול, ושקולה לתכונה \textenglish{$\oint \hvec{F}\cdot d\hvec{r} = 0$}.
  \item \textbf{אנרגיה פוטנציאלית:} \textenglish{$U(\hvec{r}_2)-U(\hvec{r}_1) = -\int_{\hvec{r}_1}^{\hvec{r}_2}\hvec{F}_C\cdot d\hvec{r}$}; כבידה: \textenglish{$E_p = mgh$}.
  \item \textbf{אנרגיה מכנית:} \textenglish{$E = K + U$}, ומשפט עבודה--אנרגיה II: \textenglish{$W_{NC} = \Delta E$}.
  \item \textbf{שימור האנרגיה:} אם אין כוחות אי-משמרים העושים עבודה --- \textenglish{$E = K + U = \text{const}$}.
  \item \textbf{סיווג:} כובד --- משמר; מתיחות/נורמל/חיכוך סטטי --- אינם עושים עבודה; חיכוך קינטי/גרר --- אי-משמרים.
  \item \textbf{כוח הגישה:} מטוטלת --- \textenglish{$v_{\max}=\sqrt{2gL(1-\cos\theta_0)}$} בשורה אחת; מכונית --- \textenglish{$v_{\max}=\sqrt[3]{P/c}$} ממאזן הספקים.
\end{enumerate}
\end{tcolorbox}

\end{document}